\begin{table}[t]
\centering
\caption{Summary Statistics: Childhood Health Outcomes by Deprivation}
\label{tab:summary}
\begin{tabular}{llrrccccc}
\toprule
 & Period & LAs & Obs & Mean & SD & Q1 & Q5 & Gap \\
\midrule
\multicolumn{9}{l}{\textit{Panel A: Dental Decay in 5-Year-Olds (\%)}} \\
 & Pre-SDIL (2007--2014) & 339 & 985 & 26.3 & 8.8 & 20.1 & 34.1 & 14.0 \\
 & Post-SDIL (2018--2023) & 324 & 794 & 22.0 & 8.0 & 15.5 & 29.3 & 13.9 \\
[0.3em]
\multicolumn{9}{l}{\textit{Panel B: Childhood Obesity, Reception Year (\%)}} \\
 & Pre-SDIL (2006--2015) & 348 & 3,279 & 22.2 & 2.9 & 19.6 & 23.9 & 4.3 \\
 & Post-SDIL (2016--2024, excl.\ 2020) & 348 & 2,644 & 22.2 & 3.1 & 19.2 & 24.4 & 5.2 \\
\bottomrule
\end{tabular}
\par\vspace{0.3em}
{\footnotesize \textit{Notes:} Dental decay is the percentage of 5-year-olds with experience of visually obvious dental decay, from the National Dental Epidemiology Programme (NDEP). Childhood obesity is the percentage of Reception-year children (age 4--5) classified as overweight or obese (BMI $\geq$ 85th centile), from the National Child Measurement Programme (NCMP). Q1 = least deprived IMD 2019 quintile; Q5 = most deprived. Gap = Q5 mean $-$ Q1 mean. COVID year 2020/21 excluded from obesity post-period due to incomplete NCMP data collection.}
\end{table}
